\documentclass[11pt,a4paper]{article}
\usepackage[T1]{fontenc}
\usepackage[utf8]{inputenc}
\usepackage{lmodern}
\usepackage{amsmath,amssymb,amsthm}
\usepackage{booktabs,longtable,array}
\usepackage[margin=27mm]{geometry}
\usepackage{microtype}
\usepackage{xcolor}
\usepackage[colorlinks=true,linkcolor=blue!45!black,citecolor=blue!45!black,urlcolor=blue!45!black]{hyperref}
\hypersetup{
 pdftitle={An upper bound for three-sunflower-free families of four-element sets},
 pdfauthor={Oleksiy Babanskyy},
 pdfsubject={Extremal set theory: three-sunflower-free four-uniform families},
 pdfkeywords={Sunflowers, four-uniform families, extremal set theory, exact certificates}
}
\usepackage{fancyhdr}
\pagestyle{fancy}
\fancyhf{}
\fancyhead[L]{\small Three-sunflower-free four-sets}
\fancyhead[R]{\small O. Babanskyy}
\fancyfoot[C]{\thepage}
\setlength{\headheight}{14pt}
\numberwithin{equation}{section}
\newtheorem{theorem}{Theorem}[section]
\newtheorem{lemma}[theorem]{Lemma}
\newtheorem{proposition}[theorem]{Proposition}
\newtheorem{corollary}[theorem]{Corollary}
\theoremstyle{remark}
\newtheorem{remark}[theorem]{Remark}
\newcommand{\F}{\mathcal F}
\renewcommand{\L}{\mathcal L}
\newcommand{\Pten}{\mathcal P}
\newcommand{\Qeight}{\mathcal Q}
\newcommand{\supp}{\operatorname{supp}}
\newcommand{\SF}{\operatorname{SF}}
\newcommand{\dd}{\mathop{}\!\mathrm d}
\newcommand{\Nclosed}[1]{N_{\F}[#1]}
\newcommand{\disj}[1]{D_{\F}(#1)}
\newcommand{\binomset}[2]{\binom{#1}{#2}}
\setlength{\emergencystretch}{2em}
\title{\vspace{-12mm}\textbf{An upper bound for three-sunflower-free families\\of four-element sets}}
\author{Oleksiy Babanskyy}
\date{}
\begin{document}
\maketitle

\begin{abstract}
Let $ M(4,3)$ be the maximum size of a finite simple family of four-element
sets containing no three distinct members whose pairwise intersections are
equal. The empty intersection is included in the forbidden configurations.
We give a computer-assisted exclusion of families of size $69$,
giving the bound $54\le M(4,3)\le 68$. The argument combines an
imported bound of $27$ for intersecting families, structural analysis of
point-links, exact trace constraints and finite exhaustions. The proof is
organized at the cardinality $69$. A cap on families containing a complete $ R16$ point-link
provides a shorter final contradiction. We make the imported results,
computational lemma contracts and remaining verification limits explicit.
The lower construction is classical. Neither exclusion of size $68$ nor
optimality of the lower construction is claimed.
\end{abstract}
\noindent\textbf{Keywords:} Sunflowers; four-uniform families; extremal set theory;
exact certificates.



\section{The finite extremal problem}
\label{sec:introduction}
A three-sunflower is a triple of distinct sets $ A,B,C $ such that
\[
 A\cap B=A\cap C=B\cap C.
\]
The common intersection is the kernel. Allowing the empty kernel means that
three pairwise disjoint sets are also forbidden. Throughout, families are
simple: a set occurs at most once.

The quantity studied here is
\[
 M(4,3)=\max\{\lvert\F\rvert:\F\text{ is a finite simple
 four-uniform family without a three-sunflower}\}.
\]
Thus an upper bound $68$ is equivalent to saying that every such family of
\emph{at least $69$} sets contains a three-sunflower. The convention that
counts the first cardinality forcing a sunflower adds one to $ M $.

\begin{theorem}[Computer-assisted bound]
\label{thm:main}
With H27 and the classification and finite-exclusion lemmas stated in this
manuscript, no simple four-uniform three-sunflower-free family of size $69$
exists. Consequently,
\[
 \boxed{54\le M(4,3)\le 68}.
\]
\end{theorem}

The general sunflower conjecture of Erd\H{o}s and Rado~\cite{ER1960}
concerns exponential dependence on the
uniformity for each fixed number of petals. A bound for this particular
four-uniform case does not resolve that asymptotic conjecture. The distinction
also prevents a finite improvement from being described as a solution of the
entire Erd\H{o}s problem.

The contemporary primary comparison is the preprint of Axante et al.
\cite{AXB2026}, whose Theorem~6.9 gives $54\le f(4,3)\le83$, with $ f $
defined there as the maximum sunflower-free size. Their Proposition~6.7
proves that an intersecting four-uniform three-sunflower-free family has at
most $27$ members, and Proposition~6.8 gives equality. We use the upper bound
as an imported theorem, denoted H27. It is not a claim that general
three-uniform families have maximum $27$.

The lower bound $54$ is attributed to Abbott and Hanson
\cite{AH1974} in the table of Axante et al. We reproduce the construction
and its elementary verification in Section~\ref{sec:lower}. No new lower
bound or priority for that construction is asserted. The numerical
comparison with $83$ provides the comparison for this bound; it does not by
itself establish that $68$ is a new record. The historical upper
bound $142$ of Abbott, Hanson and Sauer~\cite{AHS1972} is also
reported in that table; the numerical attribution here is through
the checked preprint rather than a new reading of the original article.



The imported intersecting-family theorem and the exact finite lemmas are
stated with their domains. The companion provides the complete named
finite verification route; the written reductions establish why these
finite domains cover every hypothetical family.

\section{Notation and the proof architecture}
\label{sec:notation}
We abbreviate a finite simple $r$-uniform three-sunflower-free family
as SF$r$; the forbidden configurations include the empty kernel.
Only points that occur in a member of $\F $ belong to its support. Write
$ d(x)=\lvert\{A\in\F:x\in A\}\rvert $ and
$ d(x,y)=\lvert\{A\in\F:\{x,y\}\subseteq A\}\rvert $. For $ R\in\F $,
\[
 N_{\F}[R]=\{A\in\F:A\cap R\ne\varnothing\},\qquad
 D_{\F}(R)=\{A\in\F:A\cap R=\varnothing\},
\]
and $\mu(\F)=\min_{R\in\F}\lvert N_{\F}[R]\rvert $.
The neighborhood is closed: it includes $ R $ itself. The point-link
\[
 L_x=\{A\setminus\{x\}:A\in\F,\ x\in A\}
\]
is a complete link of the ambient family. A point degree in $ L_x $ is the
global codegree $ d(x,y)$, not the global degree $ d(y)$.

\begin{lemma}[The neighborhood lower bound]
\label{lem:h27}
If $\F $ is three-sunflower-free, then $ D_{\F}(R)$ is intersecting for every
$ R\in\F $. Under H27,
\[
 \lvert N_{\F}[R]\rvert\ge\lvert\F\rvert-27.
\]
In particular, a hypothetical family of size $69$ has $\mu\ge 42$.
\end{lemma}
\begin{proof}
If two members of $ D_{\F}(R)$ were disjoint, they and $ R $ would be a
sunflower with empty kernel. H27 applies to this intersecting subfamily,
which is still four-uniform, simple and sunflower-free.
\end{proof}

The disjointness graph has the \emph{members} of $\F $ as vertices; two are
adjacent exactly when they are disjoint. This graph is triangle-free.
Its degree at $ R $ is $\lvert\F\rvert-\lvert N_{\F}[R]\rvert $.
Homomorphism fibers in this graph consist of members of $\F $, whereas the
degrees and codegrees above count points. Maintaining this distinction is
essential to the incidence arguments.

Our imported graph theorem is due to Jin~\cite{JIN1993}, in the exact
form stated as Theorem~A by Nikiforov~\cite{NIK2010}: for
$1\le k\le 9$, every triangle-free graph of order $n$ with
minimum degree strictly greater than $(k+1)n/(3k+2)$ is homomorphic
to the Andr\'asfai graph $A_k$. Its cyclic definition is given in
Section~\ref{sec:fibrebridge}. We verify the strict inequality at
$k=3$ or $k=4$ in each use; equality does not suffice.

The proof follows three reductions. First, universal link bounds and
cardinality-$69$ trace constraints exclude maximum point degrees
$20,19,18,17$ and at most $15$. In the remaining maximum-$16$ case, a
triangle-free graph homomorphism and a near-transversal argument force
every degree-$16$ point to have a complete $ R16$ link. Finally, the
local $ R61$ lemma contradicts $\lvert\F\rvert=69$. The first
reduction and all finite coverings are recorded with their hypotheses;
none is transferred from cardinality $70$ without its own bridge.

\section{Elementary link structure}
\label{sec:foundations}
\subsection{Graphs without three-sunflowers}
\begin{lemma}
\label{lem:sf2}
A simple graph contains no three-sunflower of edges if and only if its
maximum degree and matching number are both at most two. It has at most
six edges. Equality gives two vertex-disjoint triangles.
\end{lemma}
\begin{proof}
The two possible kernels of a three-sunflower of distinct edges are a
single vertex and the empty set. They respectively give a three-edge
star and a three-edge matching. Every nontrivial component of a graph
with maximum degree two is a path or a cycle. With matching number $ q $,
a path has at most $2q $ edges; an even cycle has $2q $ and an odd cycle
$2q+1$. The total matching number is at most two, so the maximum six
edges is attained precisely by two triangles. With five edges the
possibilities are a five-cycle or a triangle and a two-edge path.
\end{proof}

\subsection{The canonical ten- and eight-member families}
\label{sec:canonical}
We abbreviate $\{0,1,2\}$ as $012$ when all labels are single digits. Put
\[
 \Pten=\{012,015,023,034,045,124,134,135,235,245\}.
\]
This intersecting family has ten triples, every point degree five and
every pair codegree two. Every point-link is a five-cycle, so it is
three-sunflower-free. Write $ P_9$ and $ P_8$ for subfamilies obtained by
deleting respectively one and two blocks of $\Pten $.

On seven distinct points let
\[
 Q(h,k;T;a,b)=\{\{h\}\cup S,\{k\}\cup S:S\in\binom{T}{2}\}
 \cup\{\{h,k,a\},\{h,k,b\}\},\qquad |T|=3.
\]
This is the eight-member family $ Q_8$. The points $ h,k $ are its hubs,
$ T $ its triad and $ a,b $ its leaves. The hubs have degree five, the triad
points degree four and the leaves degree one. Pair codegrees are at
most two, and every point-link has matching number at most two.

\begin{proposition}[Intersecting triple classification]
\label{prop:i810}
An intersecting three-uniform three-sunflower-free family has at most
ten members. At ten it is $\Pten $; at nine it is $ P_9$; at eight it is
$ P_8$ or $ Q_8$, up to relabeling.
\end{proposition}
\begin{proof}
Every point-link is a graph from Lemma~\ref{lem:sf2}, so its point
degrees are at most six. A degree-six point has two triangles as its
link. A triple avoiding it would have to cover both triangles with
only three points, which is impossible; the whole family then has
only six members. A degree-five point has link either a five-cycle
or a triangle and a two-edge path. In the first case every triple
avoiding the point is one of the five three-point vertex covers of
the cycle, giving a subfamily of $\Pten $. In the second case there
are at most three such triples, and at size eight all three occur,
giving $ Q_8$.

We justify the existence of a degree-five point at sizes at least
eight. For an intersecting family $\mathcal H $ of $ m $ triples,
every $ A\in\mathcal H $ satisfies
\[
 \sum_{x\in A}d_{\mathcal H}(x)\ge m+2.
\]
If all degrees are at most four this gives $ m\le 10$; at $ m=10$ all
used points would have degree four, contrary to a degree sum of $30$.
At $ m=9$ all used points have degree three or four. If their number
is $ n $, and $ B $ counts pairs of points of codegree two, counting
intersections of pairs of triples gives $ B=45-6n $. The degree sum
$27$ and $ B\ge 0$ force $ n=7,B=3$. The triples would use
$27-B=24$ distinct pairs on seven points, exceeding $\binom 72=21$.

At $ m=8$ and maximum degree four all used degrees are at least two.
Let $ a,b,c $ count degrees two, three, four. Then $ a+b+c=n $ and
$2a+3b+4c=24$. Intersection counting gives $ B=20-3n+c $.
For each point the number $ r(x)$ of codegree-two neighbors is at
least $ d(x)-2$: the line graph of its point-link has independence
number at most two, whereas deleting at most one endpoint per edge
gives an independent set of size at least $ d(x)-r(x)$. Consequently
$2B\ge b+2c $, so $40-6n\ge b $ and $ n\le 6$. The degree sum forces
$ n=6$ and all degrees four. Now $ B=8$, so the eight triples use
$24-B=16$ distinct pairs on six points, exceeding $15$. This final
contradiction completes the classification without a support cutoff.
\end{proof}

\subsection{Isolation and a separated neighborhood}
\begin{lemma}[Isolation]
\label{lem:isolation}
In a sunflower-free triple family containing $ P_{10},P_9$ or $ P_8$,
every additional triple meeting the six-point support is one of the
deleted blocks. For $ Q(h,k;T;a,b)$, every additional triple meeting
the five-point core $\{h,k\}\cup T $ is $\{h,a,b\}$ or $\{k,a,b\}$.
\end{lemma}
\begin{proof}
For $ P_{10}$ all pairs on its support already have codegree two.
Meeting the support in one point gives a third disjoint edge in
that point-link, using a matching of size two in a five-cycle.
For $ P_9$ or $ P_8$, a new internal triple can use only pairs of the
deleted blocks. The only triangles in the union of those one or
two triangles are the deleted blocks themselves. A single-point
intersection is still excluded by a matching of size two in a
five-cycle with at most two edges removed. For a two-point
intersection $ ij $, choose an endpoint contained in just one deleted
block if the pair has remaining codegree one. Its link is a
four-edge path with the other endpoint at an end; after removing
that endpoint a matching of size two remains. If the pair has
remaining codegree zero, the two deleted blocks both contain it;
removing their two incident edges again leaves a three-edge path
with matching number two. These matchings create a forbidden
singleton-kernel sunflower with the new triple.

All ten pairs of the $ Q_8$ core have codegree two. A new triple
cannot use two core points. A single triad point has a matching of
size two within the core, which excludes it. For a single hub,
its link consists of the triad triangle and the two edges joining
the other hub to the leaves. Unless both leaves occur in the new
triple, one of those two edges and an edge of the triangle avoid
its remaining points and form a forbidden matching. This leaves
only the two stated triples.
\end{proof}

\begin{lemma}[B4 separation]
\label{lem:b4}
Let $ A $ belong to a sunflower-free triple family $\L $. If at least
eight members of $\L $ are disjoint from $ A $, its closed neighborhood
$ N_{\L}[A]$ is intersecting and has at most ten members.
\end{lemma}
\begin{proof}
The disjoint members form an intersecting family $\mathcal H $.
Proposition~\ref{prop:i810} classifies it. For a $ P $-type, isolation
forces every member of the remaining part to avoid its support:
deleted blocks cannot occur there, since they are also disjoint
from $ A $ and would belong to $\mathcal H $. Two disjoint remaining
members and a member of $\mathcal H $ would be an empty-kernel
sunflower. For $ Q_8$, the same argument uses the five-point core;
the two exceptional triples of Lemma~\ref{lem:isolation} are
disjoint from $ A $ and cannot belong to the remaining part. A
core-only block of $ Q_8$ gives the required empty-kernel sunflower
if two remaining members were disjoint. The intersecting bound
of ten now applies.
\end{proof}

The $ Q_8$ conclusion uses only core isolation. Its leaves may be
shared by the remaining part; treating all seven support points
as isolated would be incorrect.

\begin{corollary}
\label{cor:triple20}
The maximum size of a simple three-uniform family without a
three-sunflower is $20$. Therefore in a four-uniform such family,
\[
 d(x)\le 20,\qquad d(x,y)\le 6,\qquad d(x,y,z)\le 2.
\]
\end{corollary}
\begin{proof}
Relative to a triple $ A $, an exact singleton class has at most
three members and an exact pair class at most one. The residues
are intersecting sunflower-free graphs or singleton families,
respectively. Thus $|N_{\L}[A]|\le 13$. If $|\L|\ge 21$, at least
eight members avoid $ A $; B4 bounds the neighborhood by ten, and
its disjoint intersecting part also has size at most ten, a
contradiction. Two copies of $\Pten $ on disjoint supports give
twenty members. The point, pair and triple links of a four-uniform
family now give the three displayed caps.
\end{proof}

\section{Exact traces and the near-transversal argument}
\label{sec:traces}
For $ R\in\F $, let $ s_x,p_{xy},t_m $ count members with exact trace
$\{x\},\{x,y\},R\setminus\{m\}$ on $ R $. Put
$ S=\sum s_x $, $ P=\sum p_{xy}$ and $ T=\sum t_m $. Then
\[
 |N_{\F}[R]|=1+S+P+T,\qquad
 \sum_{x\in R}d(x)=4+S+2P+3T.
\]
Exact residues are intersecting and sunflower-free: two disjoint
residues with $ R $ give a sunflower with their exact trace as
kernel. In particular $ s_x\le 10,p_{xy}\le 3,t_m\le 1$.
For $ r_x=\sum_{y\ne x}p_{xy}$, the part of $ L_x $ meeting
$ R\setminus\{x\}$ has size
\[
 k_x=1+r_x+T-t_x,\qquad d(x)=s_x+k_x.
\]
When $ s_x\ge 8$, B4 makes this part intersecting. The trace
constraints keep the actual external points shared by its members;
separate integer histograms do not certify their compatibility.

Call a four-point set a \emph{complete transversal} of a triple
family if it meets every member. Relative to $ A\in\L $, a
\emph{near-transversal} is a four-point set avoiding $ A $ and
meeting every other member of $\L $.

\begin{lemma}[A near-transversal is unique]
\label{lem:near}
Let $\L $ be an SF3 family with $16$ members and, for some
$ A_0\in\L $, a B4 separation $ H=D_{\L}(A_0)$,
$ K=N_{\L}[A_0]$ into two intersecting parts of size eight.
If it has a complete
four-point transversal, each $ A\in\L $ has at most one
near-transversal.
\end{lemma}
\begin{proof}
The parts are $ P_8$ or $ Q_8$. If a part is $ P_8$, its support
is isolated from the other part. If $ H $ is $ P_8$, its deleted blocks are disjoint from
$ A_0$ and cannot belong to $ K $. If $ K $ is $ P_8$ and one member
of $ H $ is a deleted block on its support, intersection of $ H $
forces all eight members of $ H $ to meet that support; isolation
would restrict them to at most two deleted blocks, a contradiction.
Meeting at least seven members
requires at least two points from its support, since its maximum
point degree is five. All four transversal points are therefore
allocated two to each part. For two $ Q_8$ parts, their five-point
cores are disjoint; only leaves may be shared. Indeed, isolation
forces $ K $ to avoid the core of $ H $: both exceptional triples avoid
$ A_0$ and would belong to $ H $. Conversely a member of $ H $ meeting
the core of $ K $ would have to be one of its two exceptional triples;
this is impossible since every member of $ H $ uses core points of
$ H $, which lie outside the entire support of $ K $.
Each core has six
core-only triples, and each point occurs in at most four of them.
Meeting at least five core-only triples in each part requires
two points in each core. Again all four points are allocated.

A $ P_8$ part with a two-point cover must have its two deleted
blocks intersecting in two points. The $60$ automorphisms of
$\Pten $ act transitively on these $15$ pairs of deleted blocks
(out of $45$). We may take
$ P^*=\Pten\setminus\{012,015\}$, whose unique complete
two-point cover is $34$. For the canonical $ Q_8$ with hubs
$0,1$, triad $2,3,4$ and leaves $5,6$, the unique cover is $01$.
The two-point near-transversals are exactly the following:
\begin{center}
\begin{tabular}{cccc}
\toprule
$ P^*$ block & Near-cover & $ Q_8$ block & Near-cover\\
\midrule
$023$ & $45$ & $023$ & $14$\\
$034$ & none & $024$ & $13$\\
$045$ & $23$ & $034$ & $12$\\
$124$ & $35$ & $123$ & $04$\\
$134$ & none & $124$ & $03$\\
$135$ & $24$ & $134$ & $02$\\
$235$ & none & $015$ & none\\
$245$ & none & $016$ & none\\
\bottomrule
\end{tabular}
\end{center}
For an omitted block in one part, the two points in the other
part form its unique complete cover; the first pair is unique
when it exists by the table. This proves the assertion with
zero, one or two shared $ Q_8$ leaves. The finite table and
automorphism orbit are checked independently by
\texttt{primary\_check.py}; the reduction above, rather than a
chosen support cutoff, makes those finite checks complete.
\end{proof}

\section{The universal sixteen-link classification}
\label{sec:r16classification}

For an SF3 family $ L $, write $ D_L(T)=\{U\in L:U\cap T=\varnothing\}$ and
$\rho(L)=\max_{T\in L}|D_L(T)|$. In this paper, ``R16'' means
$|L|=16$ and $\rho(L)\le 7$, together with the SF3 condition. It is not an
assertion that all SF3 families of size 16 are impossible.

\begin{theorem}[Universal R16 classification]\label{thm:r16class}
Every R16 is isomorphic to the following family. Take four distinct hubs
$ K=\{u,v,a,b\}$ and two disjoint three-element sets $ D,E $, disjoint also
from $ K $. Then
\[
 R=\{i\cup P:i\in\{u,v\},\ P\in\tbinom D2\}
 \cup\tbinom K3
 \cup\{i\cup Q:i\in\{a,b\},\ Q\in\tbinom E2\}.
\]
Here and below, a point in a union denotes its singleton. The proof uses
the explicitly certified finite computations described below.
\end{theorem}

\begin{proof}
We first classify the possible seven-member intersecting family that will
occur as $ D_L(T)$. Its support has at most seven points. Here is an
arbitrary-support proof of this reduction. A degree-six point would force
only six triples: a triple avoiding it cannot cover two disjoint triangles.
For maximum degree at most four and no degree-one point, let $ b,c,d $ count
points of degrees 2,3,4. Then $2b+3c+4d=21$. If $ B $ counts used pairs of
codegree two, counting intersections gives $ B=b+3c+6d-21$.
For a point $ w $, its SF2 link has matching number at most two; if $ r(w)$
counts its codegree-two neighbours, $ r(w)\ge d(w)-2$. Hence
$2B\ge c+2d $, or $2d\ge 2b+c $. Together with
$21=3n-b+d $, this implies $ n\le 7$. Equality would force $ b=d,c=0$ and
$ n=2d $, impossible for $ n=7$. Thus this case has support at most six.

If $\ell uv $ contains a degree-one point $\ell $, the other six triples
meet $\{u,v\}$. Their distributions between only $ u $, only $ v $, and both
are $(3,3,0)$ or $(3,2,1)$, after interchanging $ u,v $. A three-edge
intersecting SF2 family is a triangle; cross-intersection confines the
two residual graphs to the same triad. The first case has support six;
the second has support seven. Finally, for a degree-five point, its
link is $ C_5$ or $ K_3\sqcup P_3$, where $ P_3$ has two edges. A triple
avoiding the point must be a three-vertex cover. In the cycle case it
lies on those five vertices; in the other case it contains two vertices
of the triangle and the middle vertex of the path. This gives support
six or the unique seven-point model, respectively.

The finite catalogue on six points is complete: every intersecting
seven-triple family selects seven of the ten complementary pairs of
triples, choosing one side from each. This gives
$\binom{10}{7}2^7=15360$ choices. The checker reconstructs all 1860 SF3
labelled families, partitions them into the seven disjoint orbits below,
and adds the analytic seven-point case. It also checks every mixed trace
of size 1 or 2 and its forbidden-pair witness, for 175 traces in total.
Strings denote sets, so $012=\{0,1,2\}$.

\begin{center}\small
\begin{tabular}{c l r l}\toprule
Type&Seven triples&Six-point orbit&Mixed traces allowed\\\midrule
$ H_0$&$012,034,134,035,235,145,245$&120&$01,02,12$\\
$ H_1$&$012,034,234,135,235,145,245$&360&$01,02,12$\\
$ H_2$&$012,134,234,135,235,145,245$&60&$0,01,02,12$\\
$ H_3$&$013,024,134,125,235,045,345$&240&none\\
$ H_4$&$023,124,134,125,235,045,345$&360&none\\
$ H_5$&$014,234,025,125,035,145,345$&360&$01,23$\\
$ H_6$&$134,234,125,035,135,045,245$&360&$0,12$\\
$ H_7$&$023,024,034,124,134,015,016$&---&$5,6,23,56$\\\bottomrule
\end{tabular}
\end{center}
For a forbidden trace $ X $, the certificate exhibits two actual triples
$ h_1,h_2$ with $ h_1\cap h_2=h_1\cap X=h_2\cap X $. Adding points outside
the catalogue support does not change these intersections. Thus this is
a trace obstruction on arbitrary support, rather than a cutoff on new
point labels. Allowed traces are not assumed jointly realizable.

We next rule out $|L|=16,\rho(L)\le 6$. Every closed triple-neighbourhood
has size at least ten. A degree-one point cannot have a degree-six
neighbour, because every used pair in the link of such a neighbour has
codegree two. Its other two neighbours therefore have degrees at most
five, and its root neighbourhood has size at most nine. For a degree-two
point $ x $ with a degree-six neighbour $ a $, its triples are $ xab,xac $,
and the link of $ a $ forces $ abc $. Hence
$ d(xb)=1,d(ab)=d(ax)=2,d(b)\le 5$, giving $|N_L[xab]|\le 9$.
Without a degree-six neighbour, a root $ xyz $ could have neighbourhood
ten only with $ d(y)=d(z)=5$ and all three used pairs of codegree one.
But the link of $ y $ then has $ x $ as a leaf of its $ P_3$ component and
$ z $ as that component's middle vertex, forcing $ d(yz)=2$.
The minimum used degree is consequently at least three.

Let $ I $ count intersecting pairs of triples. Then $ I\ge 72$ and
\[
 I=\sum_w\left(\binom{d(w)}2-\frac{r(w)}2\right),\qquad \sum_w d(w)=48.
\]
At degrees 3,4,5,6 the minimum values of $ r(w)$ are 1,2,4,6.
For those degrees the summand is at most $19d(w)/6-7$, so
$72\le 152-7n $ and $ n\le 11$. If the maximum degree is at most five,
the stronger upper bound $(11d(w)-23)/4$ gives $ n\le 10$; a degree-five
point then exists, because otherwise $48\le 4n\le 40$.
The following three padded domains therefore cover every remaining
family. Padding permits unused points and does not discard smaller
supports.

\begin{center}
\begin{tabular}{l r r r}\toprule
Complete normalized point-star&Universe&Proof-tree nodes&Solutions\\\midrule
Two triangles&11&16357&0\\
$ C_5$&10&5874&0\\
Triangle and two-edge path&10&3641&0\\\bottomrule
\end{tabular}
\end{center}
For the first domain, six triples through point 0 are fixed and ten of
the 120 triples avoiding 0 are required. In the other domains five are
fixed and eleven of the 84 triples avoiding 0 are required. The supplied
proof trees cover every choice. Their set-based checker verifies each
branch and each terminal capacity, degree, disjointness or colouring
bound. In particular, all 99953 pairs in a colour class conflict with an
actual triple of the fixed star. A branch is assigned by its last chosen
candidate in the certified order; the omitted prefix has fewer colours
than the number of triples still required. The three supplied complete
trees have 25872 nodes in total and no open
leaves. The files are \path{companion/terminal/r16/results/r16_mu10_certificate.json.gz}
and \path{companion/terminal/r16/src/check_r16.py}. Thus $\rho(L)\le6$ is impossible.

Choose $ A\in L $ with exactly seven disjoint triples. Put $ H=D_L(A)$,
$ W=\operatorname{supp}H $ and $ K'=L\setminus H $. The family $ H $ is
intersecting and has one of the eight types above. No member of $ K'$ is
entirely in $ W $, since it would be disjoint from $ A $ and belong to $ H $.
The part $ K_0$ disjoint from $ W $ is intersecting, and $|K_0|\le 7$ because
each member of $ H $ avoids all of it.
For $ H_3,H_4$ there are no mixed members, giving $|K_0|=9$.
For $ H_0,H_1,H_5$ there are at most three mixed members, each on a pair
already used once. Hence there are at least two mixed members and at
least six in $ K_0$. If a mixed member is $ P+r $, take a member of $ H $
avoiding $ P $. Every member of $ K_0$ must contain $ r $ to avoid three
disjoint triples. Distinct external points for two mixed members would
force a common pair in at least six members of $ K_0$; a single common
external point instead has degree at least eight. Both are impossible.
For $ H_6$, its block 125 is disjoint from all of $ K'$ except possibly the
single member with trace 12, giving at least eight disjoint members.
The same argument uses 023 and 23 for $ H_7$.

Only $ H_2$ remains. Its six-triple core can be renamed
$\{u+P,v+P:P\in\binom D2\}$, and its seventh member is $ uvb $.
Pair capacities and the two triples $ uij,vik $ forbid any further triple
meeting $ D $. All remaining members have the forms $ uvw $, $ u+e $, $ v+f $,
or avoid $\{u,v\}\cup D $; call the last part $ K^*$. The part $ K^*$ is
intersecting and has size at most seven. There are either one or two
members of the first form, by the capacity of $ uv $.
If $ uvb $ is the only such member, the other two forms can only be
$ ubr,vbs $, at most one each. Size 16 would require both and $|K^*|=7$.
For $ r\ne s $, members of $ K^*$ avoiding $ b $ contain $ rs $, at most two;
members containing $ b $ have residual edges meeting $ r $ or $ s $, each
with at most one remaining incidence, at most two more. For $ r=s $,
the pair $ br $ is already saturated and $ K^*$ splits into the $ b $ and
$ r $ stars. Their residual graphs are cross-intersecting. If the first
has 0,1,2,3 edges, the second has at most 6,4,3,3, respectively. The
two-edge case is a two-edge path and the three-edge case a triangle.
Thus $|K^*|\le 6$ in every case, again impossible.

With two members $ uva,uvb $, each extra edge through $ u $ or $ v $ must be
$ ab $, so each contributes at most one triple. At least one is necessary
for size 16. Every member of $ K^*$ then meets $\{a,b\}$. In each hub-link
the edge $ uv $ is already present, so residual edges form an intersecting
SF2 family, at most three edges. Therefore $|K^*|\le 6$. Equality in
$6+2+1+1+6=16$ forces both extra triples and two residual triangles,
one through $ a $ and one through $ b $. Cross-intersection forces the same
triad $ E $, disjoint from all existing core points. This is precisely $ R $.
\end{proof}

\begin{lemma}[Maximality and the high-link decompositions]\label{lem:r16maximal}
No new triple, including one with new points, can be added to $ R $.
Consequently every SF3 family of 17 triples is $\Pten\sqcup I_7$ or
$ P_9\sqcup I_8$, and every such family of 18 triples is
$\Pten\sqcup I_8$ or $ P_9\sqcup P_9$, on disjoint supports.
Here $ I_j $ denotes an intersecting SF3 family of size $ j $, and
$ P_9$ is $\Pten $ with one block removed.
The size 19 and 20 decompositions are respectively
$\Pten\sqcup P_9$ and $\Pten\sqcup\Pten $, again on disjoint supports.
\end{lemma}
\begin{proof}
Every hub has degree six and its complete SF2 link is already maximum.
A new triple avoids all hubs. Used pairs within $ D $ and within $ E $ have
codegree two, so it uses at most one point in each triad. A pair of $ D $
and a pair of $ E $ avoiding the new triple give two disjoint members of
$ R $, disjoint also from that triple. This is a forbidden empty sunflower.
If a family of 17--20 triples had at most seven disjoint triples per
root, every sixteen-member subfamily would be $ R $ and could not extend.
Thus some root has at least eight disjoint triples. B4 separates the
family into two intersecting parts of size at most ten. Their size pairs
are 10+7 or 9+8 at size 17,10+8 or 9+9 at size 18,10+9 at size 19,
and 10+10 at size 20. Isolation makes their
supports disjoint: a missing block of a P9 belongs to the same
disjointness class as that P9 and cannot be inserted in the other part.
\end{proof}

\section{The direct reduction of a hypothetical family of size 69}
\label{sec:f69branches}

This section gives the local finite exclusions at cardinality $69$.
The companion includes their full mathematical reductions, coefficient
constructors, exact certificates and coverage checkers. The nine blocks
are identified by their mathematical roles in \path{proof/DEPENDENCIES.md}.
All file paths in this manuscript are relative to the companion repository root.
All domains below use $\mu\ge42$; no exclusion at a larger cardinality is
used as a premise.

\begin{proposition}[Bootstrap at 42]\label{prop:bootstrap42}
In an SF4 family with $\mu\ge 42$, a singleton of size 9 or 10 at an actual
anchor requires at least seven other members through its centre.
Every used point has degree at least four; co-occurring points have
degree sum at least 20, and every member has degree sum at least 52.
Moreover
\[
 \frac{37}{170}|\F|\le |\operatorname{supp}\F|
 \le\frac{17}{50}|\F|.
\]
\end{proposition}
\begin{proof}
For an ordered anchor $ R=\{r_0,r_1,r_2,r_3\}$ let $ p_{ij}$ be its six
exact pair traces, $ t_i $ its exact triple trace omitting $ r_i $, and $ s_i $
its singleton traces. Then
\[
 k_i=1+\sum_{j\ne i}p_{ij}+\sum_h t_h-t_i,\quad d(r_i)=k_i+s_i,\quad
 |N[R]|=1+\sum p+\sum t+\sum s.
\]
The intervals are $0\le p\le 3$, $ t\in\{0,1\}$, $0\le s\le 10$.
The universal rooted signatures of the intersecting triple models and
the high-link decompositions are imposed only in their valid cases.
The full extension-point identifications, all five trace margins and the
high-point bound are proved in
\path{companion/upstream/MU42/TRACE_FOUNDATIONS.md}.
The two complete row constructors yield 980624 degree-ordered rows and
1233 degree profiles. They imply minimum degree 2 and minimum co-occurring
degree sum 20. Permutations order an actual anchor; they do not impose
symmetry on its completion.

For an isolated P10 or P9 on six points, all internal quadruples are
chosen, including absences. The respective catalogues have 292 and 373
members, covering all $2^{15}$ masks. At each of the six neighbour
points, the actual external triple family must be shared by every
central anchor that mentions it. These families are cross-intersecting,
and a triple appears in at most two of them. If one has size 10 their
total is at most 20; if one has size 9 and none has 10 it is at most 22,
by the complete triple-transversal catalogues. A singleton 8 in a Q8
only guarantees its six core triples are external; it is not replaced
by an external family of size 8.
The 256 correlated lower-bound vectors in $\{0,6,9,10\}^6$ cover every
actual choice. The exact joint check eliminates all 16 P10 and 24 P9
scalar survivors with $ k\le 5$.

At $ k=6$, the full 13 P10 interior orbits are checked before the 79 P9
orbits. The P9 check may therefore use the already proved P10 bound 7.
Each system uses the 41 nonempty traces of sizes 1--3 on the six points;
empty traces and the number of exterior points are unrestricted.
The same variables encode the same members in all central anchors.
Ninety systems have an empty anchor domain. The two remaining P9 systems
have exact integer certificate gaps $-53$ and $-200$, including all
variable-box contributions. Thus $ k\ge 7$ for both singleton sizes.

The only degree-two raw row is
$ p=(0,0,0,3,3,3)$, $ t=(1,0,0,0)$,
$ s=(1,10,10,10)$, $ d=(2,18,18,18)$.
Thus a degree-two point $ x $ has two members $ xA,xB $ with disjoint triples
$ A,B $, whose six points all have degree 18. In every $ L_a $, $ a\in A $,
the point $ x $ is a leaf, so the 18-link decomposition forces P10+Q8,
with $ x $ a Q8 leaf and $ A\setminus\{a\}$ its two hubs. The other leaf
is the same point $ y $ for all $ a $: otherwise the global triple $ A $
would have more than two completions. Comparing the actual members
through a pair $ a,b\in A $ forces the same Q8 triad $ C $ in all three
links. Hence $ xA,yA $ and all nine members
$ E_A\cup E_C $, $ E_A\in\binom A2,E_C\in\binom C2$, are present.
Each opposite P10$ H_a $ avoids $ A\cup C\cup\{x,y\}$.
The point $ y $ is not in $ B $: otherwise its degree 18 link would contain
$ A $ with all three points of degree one in that link, which no 18-link
model permits.

If $ B\cap C=\varnothing $, then $ B $ avoids all of that core. For
$ T\in H_a $, a disjointness $ B\cap T=\varnothing $ would make $ xB,aT $
and $(A\setminus\{a\})\cup E_C $ three disjoint members. Thus $ B $
is a triple-transversal of every $ H_a $, and therefore a block of each
P10. The three actual members $ aB $, $ a\in A $, form a sunflower with
kernel $ B $.
If $ c\in B\cap C $, the six prescribed triples in $ L_c $ are
$ t\cup E_A $, $ t\in C\setminus\{c\}$ and $ E_A\in\binom A2$.
They occupy the same intersecting 18-link component. The point-links
of all three points of $ A $ contain $ C_4$, whereas every point-link of
P10 is $ C_5$ and contains no $ C_4$ subgraph. Thus that component is
Q8 with triad $ A $ and hubs $ C\setminus\{c\}$. Its leaf $ x $ forces
$ B\setminus\{c\}=C\setminus\{c\}$, so $ B=C $.
For $ a\in A,c\in C $, disjoint tails $ T\in H_a,V\in H_c $ would make
$ aT,cV,(A\setminus\{a\})\cup(C\setminus\{c\})$ three disjoint
members. These P10s are cross-intersecting, so each is contained in
the triple-transversal catalogue of the other and they coincide.
A shared block again supplies three members $ aT $ with kernel $ T $.
This excludes the degree-two row on arbitrary support.
Each of the four degree-three raw rows has singleton 2 at that point
for every one of its roots. The three triples of its complete link
would consequently be pairwise disjoint, a forbidden sunflower at
the point. There is no degree-one raw row. Reconstructing the filtered
domain gives 292560 rows and 966 degree profiles, with member degree
sum at least 52.

For completeness, the support calculation is a global identity rather
than a maximization of unrelated anchors. Put $ c(R)=\sum_{v\in R}1/d(v)$.
The verified pointwise row inequality is
\begin{align*}
10000c(R)\le 3400+26G_4(R)-33B_{17a}(R)-21B_{17b}(R)
 -38B_{18}(R)+9B_{19}(R)+18T_{16}(R).
\end{align*}
Here $ G_4=\sum_{v\in R,d(v)=4}(4s_v-8)$; the high-link corrections are
\begin{align*}
B_{17a}&=\sum_{d(v)=17}(7[s_v=7]-10[s_v=10]),&
B_{17b}&=\sum_{d(v)=17}(8[s_v=8]-9[s_v=9]),\\
B_{18}&=\sum_{d(v)=18}(8[s_v=8]-10[s_v=10]),&
B_{19}&=\sum_{d(v)=19}(9[s_v=9]-10[s_v=10]),\\
T_{16}&=\sum_{d(v)=16}(9[s_v=9]-7[s_v=7]).
\end{align*}
All displayed sums are over the indicated points of the same $ R $.
The $ B $ corrections sum to zero across all members, by their complete
two-part link decompositions; $\sum_R G_4(R)\le 0$ follows from the
triangle-free disjointness graph of a four-triple point-link and
Mantel's elementary bound, and $\sum_R T_{16}(R)\le 0$ follows from
B4 and isolation of any P9 in a sixteen-link. The positive coefficients
26 and 18 preserve these inequalities after summation. The row check
is performed after multiplying by $\operatorname{lcm}(1,\ldots,20)$;
its exact margin is 181/85000. Since $\sum_R c(R)=|\operatorname{supp}\F|$,
the upper support bound follows. The pointwise lower bound 37/170 gives
the other bound. The source, finite objects and exact arithmetic contract
are listed in \path{proof/DEPENDENCIES.md}, under
\emph{MU42: Neighbourhood bootstrap and separate C5}.
\end{proof}

\begin{proposition}[Complete F69 degree reduction]\label{prop:f69reduction}
A hypothetical SF4 family of 69 members satisfying H27 has maximum degree 16.
Every exact singleton has size at most 8 and every closed neighbourhood
has size at most 44.
\end{proposition}
\begin{proof}
H27 supplies $\mu\ge 42$ directly. The previous proposition gives minimum
used degree 4 and support 16--23. The classified maximum-degree cases and
their finite covers are as follows. All quantities refer to necessary
local domains, not to counts of realized global families.
\begin{center}\small
\begin{tabular}{c p{9.5cm}}\toprule
Maximum degree&Exclusion and final ordinary contradiction\\\midrule
20&P10+P10. One internal certificate and six external certificates force
each residual member to use at least two exterior points. A mixed-member
cover then has capacity 44, below the required 49. Without a mixed member,
each of 49 residual members avoids at least 12 of the 20 cones, giving
$12\cdot 49=588>20\cdot 27=540$.\\
19&P10+P9. All P9 interiors are eliminated before empty P10 is imposed.
The mixed cover retains the pure-three P5 exceptions and has capacity 44,
below 50. The terminal three-unit deficit gives $547>19\cdot 27=513$.\\
18&All P10+I8 roots are covered by 12 internal certificates and 18 complete
extension forests. The P9+P9 root uses the same external identities in
both trace systems:20 internal,22 mixed-four and 8 mixed-three exact
certificates. Mixed-two has capacity 44 below 51. The joint external
capacity of the two complementary P9 traces is 2, giving deficit 22 and
$488>18\cdot 27=486$.\\
17&The P10+I7 and P9+I8 alternatives are both covered. The P10 bridge
regenerates the 17/42 domain and identifies all 18249 states with two
universal finite cores; it imports no F70 family. Twelve P10 and 75 P9
certificates, together with 870 complete extension forests, exclude 17.\\
$\le 15$&All 28 high-neighbourhood row orbits and 35 realized cores are covered.
Thirty-three fail another actual anchor; two equivalent remaining cores
share one 87-trace system. Its exact certificate has $ H=-435,L=-89$,
contradicting $ L\le H $. Thus all closed neighbourhoods are at most 43,
contradicting the graph lemma below.\\\bottomrule
\end{tabular}
\end{center}
Every LP certificate excludes a relaxation containing every actual
family in its stated root. For $ Az\le b,Ez=f,l\le z\le u $, an integer
certificate $ y\ge 0,\eta $ gives $ q=A^Ty+E^T\eta $ and
\[
 L=\sum_j\min(l_jq_j,u_jq_j)\le q\cdot z
 \le H=y\cdot b+\eta\cdot f.
\]
The decisive comparison is $ H<L $. Forest verifiers independently
reconstruct all children and the reason for each closed leaf. The
source-to-model maps keep actual anchors, shared degrees and shared
exterior identities; they do not assume that extensions symmetric in
the labels exist. The exact source sections, branch counts, files and
companion checks are recorded in \path{proof/DEPENDENCIES.md}, under
\emph{Preliminary finite blocks}.

We justify the graph contradiction used in the last row. The
disjointness graph $ G $ has vertices the members of $\F $, is
triangle-free, and has $\alpha(G)\le 27$. If all neighbourhoods have
size at most 43, its minimum degree is at least 26. Jin with $ k=3$
applies because $26\cdot 11>4\cdot 69$. The cyclic $ A_3$ definition
below is carried to Wagner labels by $ i\mapsto 3i\pmod 8$:
cyclic differences $3,4,5$ become $1,4,7$. In the Wagner labelling of
$ A_3$, let the fibre sizes be $ a_i $ and put
$ b_i=a_{i-1}+a_{i+1}+a_{i+4}\le 27$. Set $ e_i=27-b_i $;
$\sum e_i=9$, and the inverse incidence identity gives
$ a_i=12-e_{i-1}-e_i-e_{i+1}\ge 3$. Thus every fibre is occupied and
$ b_i\ge 26$. But $8\cdot 26\le\sum b_i=3\cdot 69=207$, impossible.

It remains to state the two finite reductions at maximum 16. The full
P9 interior cover has 373 assignments in 79 orbits; the LOW15 bonus
excludes 77 and leaves only the empty assignment and the three labelled
copies of 0345. NO9 closes these with 60 newly checked forests and 22777
closed leaves, without importing the old 17-degree domain. Consequently
no singleton 9 remains; singleton 10 already requires $ k\ge 7$ and cannot
occur at degree 16. The NO45 cover has 36 high-anchor orbits and 2335
actual cores. Of these,2330 fail another anchor and five exact systems
exclude the rest, so no closed neighbourhood reaches 45. This gives
the claimed singleton and neighbourhood bounds.
\end{proof}

\section{The fibre bridge forces an R16 point-link}
\label{sec:fibrebridge}

We use the cyclic definition $ V(A_k)=\mathbb Z/(3k-1)$ and
$ i\sim j $ when $ k\le j-i\le 2k-1$ modulo $3k-1$. Jin provides an
image $ A_4$ because the disjointness graph now has minimum degree 25 and
$25\cdot 14>5\cdot 69$. Deleting an unoccupied fibre and applying the
explicit edge-preserving deletion map $ A_k-v\to A_{k-1}$ yields a
surjective image with $ k\le 4$. The cases $ k=1,2$ are impossible: the
maximum independent-fibre mass would give respectively $69\le 2\cdot 27$
or $2\cdot 69\le 5\cdot 27$. Thus $ k=3$ or 4.

Here are the required deletion maps for $ v=0$; entries list the images
of vertices $1,\ldots,3k-2$ in their cyclic order:
\[
\begin{array}{c|l}
k& A_k-0\longrightarrow A_{k-1}\\\hline
2&(0,0,1,1)\\
3&(0,1,1,2,3,3,4)\\
4&(0,7,6,6,5,4,3,3,2,1).
\end{array}
\]
Checking the cyclic edge differences verifies every edge. Rotation
gives the map for any deleted vertex. If the resulting image has
another empty fibre, repeat; $ k $ strictly decreases. The displayed maps supply the entire deletion argument without an
additional archived witness.

Let $ a_i $ be the fibre masses and $ b_i=\sum_{j\sim i}a_j $.
Neighbourhoods of vertices of $ A_k $ are independent, so $ b_i\le 27$.
Surjectivity and the minimum degree give $ b_i\ge 25$. With
$ r_i=b_i-25$, the inverse fibre identity is
\[
 r_i\in\{0,1,2\},\qquad \sum_i r_i=25-6k,\qquad
 a_i=6+r_i+r_{i+k}+r_{i+2k-1}.
\]
The formula is an identity for this cyclic incidence matrix; it does
not assume that all permissible edges between fibres are present.

\begin{proposition}[Every degree16 point has an R16 link]\label{prop:allr16}
Under the conclusions of Proposition~\ref{prop:f69reduction}, every
point of degree 16 in $\F $ has an R16 complete point-link.
\end{proposition}
\begin{proof}
Fix such a point $ x $, and let $ t_i $ count its star members in fibre $ i $.
Set $ I=\{i:t_i>0\}$ and $ J=\{j:N_{A_k}(j)\cap I=\varnothing\}$.
Every actual member of a fibre in $ J $ avoiding $ x $ meets all star
members of $ x $, hence is a full four-transversal of $ L_x $. These members
are distinct, and their number is
$ B(x)=\sum_{j\in J}(a_j-t_j)$.
Because two star members cannot be adjacent in the disjointness graph,
every occupied star fibre satisfies
\[
 \sum_i t_i=16,\qquad \sum_{j\sim i}t_j\le r_i\quad(i\in I).
\]
The second inequality counts absent permissible edges of any one star
member; its budget is $ b_i-25=r_i $, which can be 2. It is not assumed
to be a matching. If $ I $ is independent, choose a cyclic window $ W $
of $ k $ consecutive indices with $ I\subseteq W $. Such a window is
independent and is the neighbourhood of a vertex of $ A_k $, so
$ I\subseteq W\subseteq J $ and $\sum_{j\in W}a_j\ge25$.
Since $\sum_{j\in J}t_j=16$, it follows that
$ B(x)=\sum_{j\in J}a_j-16\ge9$.
The remaining finite domain consists of all positive star vectors on
nonindependent supports subject to the displayed inequalities and
$ t_i\le a_i $. A nonisolated coordinate has $ t_i\le 2$, since an
adjacent occupied coordinate has budget at most 2; isolated coordinates
retain their entire possible interval. No point-star vector is lost.

The complete census and its new set-based reconstruction give:
\begin{center}
\begin{tabular}{c r r r}\toprule
Image&Dihedral fibre profiles&Nonindependent star vectors&$ B=1$ vectors\\\midrule
$ A_3$&70&1293&52\\
$ A_4$&1&0&0\\\bottomrule
\end{tabular}
\end{center}
All vectors have $ B\ge 1$. For every one of the 52 exceptional vectors,
in 23 profiles, a certified index $ i $ satisfies
\[
 \sum_{j\sim i}t_j=1,\qquad a_i-t_i\ge 9.
\]
This census is reconstructed by \path{companion/terminal/check_f69_bridge.py}; its
output lists all 52 witness vectors and indices.

Suppose $ L_x $ is not R16. A root has at least eight disjoint triples;
by the singleton bound these are exactly eight, and B4 gives two
intersecting parts of size 8. The separated parts are P8 or Q8, with
disjoint supports whenever a P8 occurs and with at most shared leaves
when both are Q8. A full four-transversal uses two points in each P8
support, or two points in each of the disjoint five-point Q8 cores.
Consequently it is unique. Only P8s whose two omissions meet in two
points have a two-cover; the 15 such P8s form one orbit, represented
by $ P^*=\Pten\setminus\{012,015\}$. The quasi-transversal lemma and
table in the preceding foundations section allow at most one
four-set avoiding a prescribed triple and meeting the other fifteen.
The complete 80-case check covers $ P^*/P^*$, $ P^*/Q8$, and $ Q8/Q8$ with
0,1,2 shared leaves; it has 54 nonempty quasi-transversal sets and
maximum cardinality 1. Arbitrary exterior points cannot occur because
the four available positions are already exhausted by the two covers.

Thus $ B(x)=1$, and $ x $ realizes one of the 52 exceptional vectors.
For its witness fibre $ i $, there is exactly one star member $ A $ in an
adjacent fibre. Every one of the at least nine members of fibre $ i $
avoiding $ x $ meets the other fifteen star members automatically, since
their fibres are nonadjacent. The unique full transversal belongs to
$ J $, whereas $ i\notin J $, so these nine members are distinct from it.
None can also meet $ A $, since that would be a second full transversal.
They therefore avoid $ A $ and are nine distinct quasi-transversals of
$ L_x $ relative to the same triple $ A\setminus\{x\}$. This contradicts
the upper bound 1. Hence $ L_x $ is R16.
\end{proof}

\section{A complete R16 point-link gives a local upper bound 61}
\label{sec:r61}

\begin{lemma}[Crowding bound]\label{lem:crowding26}
In an SF4 family with $\mu\ge 41$ and $\Delta\le 18$, the three largest
exact singleton counts of any one actual anchor sum to at most 26.
\end{lemma}
\begin{proof}
Each singleton tail is intersecting SF3 and has size at most 10; all
anchor centres avoid every tail support. A violation contains actual
three-colour subfamilies of sizes $(10,10,7),(10,9,8)$ or $(9,9,9)$.
They are rainbow-sunflower-free, since their lifts have distinct
centres. Normalize the first P10/P9 on six points. Every second support
consists of a subset of these six and consecutively renamed fresh
points. Testing every labelled P10 and omission gives 768,7680,7680
positions, in 18,59,411 orbits, respectively. Only automorphisms of
the first colour and fresh-point permutations are factored out.

For a third tail, its intersection with the first two supports determines
all rainbow intersections. Every trace of size 0--3 is tested against
all pairs from the first two colours. In the $(10,10,7)$ regime,15 of
the 18 first-pair orbits have an actual anchor with certified
neighbourhood at most 40. The cover capacities subtract the twenty
already present members from global point/pair/triple/four-set
capacities 18,6,2,1. In two of the remaining orbits all permitted tails
are internal and no seven-subset is intersecting. In the last they
are blocks on disjoint P10 supports and an intersecting third family
must lie in one support, yielding cross-intersecting stars of 10+7.

For $(10,9,8)$, one first-pair orbit is already cross-intersecting, two
have neighbourhood covers 33, and the remaining 56 admit exactly 90
third P8/Q8 records, each containing an actual cross-intersecting pair.
For $(9,9,9)$, three first-pair orbits are cross-intersecting; the
other 408 admit exactly 20 third P9 records, again with such a pair.
All support identifications, including fresh Q8 leaves, are retained.

The cross-star theorem used here has the following precise guard: both
centres are outside both tail supports, both tail families are
intersecting SF3, the two tail families are cross-intersecting, and all
lifted members are actual members of the
same SF4. For total size at least 16, there is an actual root with
$|N[R]|\le d(x)+d(y)-d(xy)+4\le 40$. Its arbitrary-support reduction
extracts 10+6,9+7, or 8+8 tails. P10/P9 triple transversals are internal;
an 8-tail has a unique two-cover or none, and at most two triples
using that two-cover with new points. Thus all 433 normalized pairs
are covered on a derived finite support. Their resource certificates
cover all 21567 useful traces; the residual capacity is 4 in 420 cases,
2 in 12 cases, and 0 in one case. The bound counts members through a
centre once as $ d(x)+d(y)-d(xy)$, and charges every remaining member
to a triple or four-set resource after subtracting its occurrences
in the sixteen known members. This proves the quoted root inequality.

All three regimes therefore contradict $\mu\ge 41$. The finite
objects and independent set-based checkers are
\path{companion/terminal/crowding/certificates/crowding26.json.gz},
\path{companion/terminal/crowding/code/verify.py},
\path{companion/terminal/cross15/results/cross16_certificate.json} and
\path{companion/terminal/cross15/src/verify.py}. The written reductions above
justify their application on arbitrary support; a finite count alone
would not do so.
\end{proof}

\begin{theorem}[R61]\label{thm:r61}
Let $\F $ be SF4 with $\mu\ge 41$ and $\Delta\le 18$. If a complete
point-link is R16, then $|\F|\le 61$.
\end{theorem}
\begin{proof}
Use Theorem~\ref{thm:r16class} to label this complete link by four hubs
$ K=\{0,1,2,3\}$, peripheral triads $ D,E $, and centre $ x $.
There are sixteen members through $ x $, including four central members
$ R_i=\{x\}\cup(K\setminus\{i\})$.
Each hub $ i $ has codegree $ d(xi)=6$. A point of global degree 17 or 18
cannot have such a neighbour: its link separates into intersecting
parts of size at least seven on disjoint supports, and each such
part has point degree at most five. Thus all four hubs have global
degrees at most 16, even though the ambient cap is 18.

Among members avoiding $ x $, use exact traces on the same $ K $: singleton
counts $ s_i $, pair counts $ p_{ij}$, triple counts $ t_i $ omitting $ i $, and
four-set count $ q $. Put
\[
 P=\sum p_{ij},\quad T=\sum t_i,\quad E_*=\sum s_i+P+T+q,
 \qquad d_i=6+s_i+\sum_{j\ne i}p_{ij}+T-t_i+q.
\]
The empty trace is unrestricted. Exact root identities and capacities
are
\begin{gather*}
 0\le s_i\le 10,\quad 0\le p_{ij}\le 3,\quad t_i,q\in\{0,1\},
 \quad t_i+q\le 1,\quad p_{ij}+t_\ell\le 3\quad(\ell\notin\{i,j\}),\\
 |N[R_i]|=16+E_*-s_i\ge 41,\qquad d_i\le 16.
\end{gather*}
For $ i\ne j $, the exact singleton of $ j $ at $ R_i $ has size
$ u_{ij}=s_j+p_{ij}$. Hence $ u_{ij}\le 10$ and
$\sum_{j\ne i}u_{ij}\le 26$ by Lemma~\ref{lem:crowding26}.
If $ u_{ij}\ge 8$, B4 makes the complementary part of $ L_j $
intersecting. It contains the six triples through $ x $; an intersecting
SF3 with a degree-six point has exactly six members. Therefore
\[
 u_{ij}\ge 8\quad\Longrightarrow\quad d_j-u_{ij}=6.
\]

If $ d_j=16$, its complete link is either 10+6 or R16, because $ x $
has degree six there. In 10+6, the opposite P10 avoids both triangles
of the six-star, giving $ u_{ij}=10$ for all $ i\ne j $. In R16 the
triangle $ K\setminus\{j\}$ is either the other-hub triangle or a
peripheral triad; all three singleton values are then 3 or 7. Its used
pairs have codegree two, so each opposite global triple has codegree
two as well. The full necessary guard is
\[
 d_j=16\Longrightarrow
 (u_{ij}:i\ne j)\in\{(3,3,3),(7,7,7),(10,10,10)\},
\]
and in the first two alternatives $ t_i+q=1$ for $ i\ne j $.
No other 16-link is lost: the separated 9+7 and 8+8 parts have point
degree at most five, with shared Q8 leaves having degree at most two.

There are $4^6\cdot 2^4\cdot 2=131072$ choices of $(p,t,q)$.
For each, enumerate all four singleton counts satisfying the exact
degree and neighbourhood inequalities. A second construction enumerates
the total $ E_*$ and its bounded compositions, using
\[
 E_*\ge\left\lceil\frac{100-P-T-q}{3}\right\rceil,
 \qquad E_*\le 40-P-2T-3q.
\]
These inequalities follow by summing the four root-neighbourhood and
degree constraints; they do not constrain exterior labels. The exact
reconstruction gives 21652 profiles before the singleton/B4 guards,
111 after them, and 11 after the complete-link guard. The eleven
records are exactly the following labelled forms.
\begin{center}\small
\setlength{\tabcolsep}{4pt}
\begin{tabular}{c l l l c l r}\toprule
Form&$ s $&$ p $&$ t $&$ q $&Hub degrees&Records\\\midrule
A&Permutations of $(8,8,9,9)$&all 0&all 0&0&$(14,14,15,15)$&6\\
B&$(7,7,7,7)$&2 on a perfect matching&all 0&0&$(15,15,15,15)$&3\\
C&$(7,7,7,7)$&all 0&all 1&0&$(16,16,16,16)$&1\\
D&$(6,6,6,6)$&all 1&all 0&1&$(16,16,16,16)$&1\\\bottomrule
\end{tabular}
\end{center}
The finite proof is \path{companion/terminal/cap61/certificates/core_certificate.json}
with its complete reconstruction in \path{companion/terminal/cap61/src/verify.py}.

In form B choose a matching edge $ ij $. The singleton of $ j $ at $ R_i $
has nine tails, hence is P9, but $ i $ occurs in exactly two of them.
All positive point degrees of P9 are 4 or 5. In form D each hub-link
is R16 with $ K\setminus\{j\}$ peripheral, whereas the actual member $ K $
would insert that forbidden peripheral triad. In form C take the
actual member $\{1,2,3,z\}$ furnished by $ t_0=1$.
Its reading in $ L_1$ forces $ z\in D $: the peripheral pair 23 can be
completed only by the centre $ x $ or its partner hub, and $ z\ne x $.
The same member read in $ L_2$ forces $ z\in E $. Since $ D\cap E=\varnothing $,
this is impossible. The outside label is kept identical in both links.

Only the six form-A records remain. For a hub $ i $ with $ s_i=9$,
$ E_*=34$ and $|N[R_i]|=16+34-9=41$ exactly. The complement $ D(R_i)$
is intersecting. Its complete $ i $-link consists of the nine pure
singleton tails, hence is exactly P9: all other nonempty hub traces
are zero, and all six members through $ xi $ are removed with $ N[R_i]$.

We prove the required ordinary bound for this complement. If an
intersecting SF4 family $\mathcal I $ has a complete point-link $ H=P9$
or $ P10$ at $ a $, put $ W=\operatorname{supp}H $ and $ h=|H|$.
Every member avoiding $ a $ covers $ H $, so uses at least three points of $ W $.
Let $ b $ count internal four-sets and $ e $ count members with a three-point
trace on $ W $. If $ h=10$, every internal four-set contains two P10
blocks, whose remaining triple capacities give $ e+2b\le 10$ and
$|\mathcal I|=10+b+e\le 20$. If $ H=P10\setminus\{T\}$,
every internal four-set contains at least one block of $ H $, so
$ e_H+b\le 9$. Its only other possible triple traces are $ T $ and
$ W\setminus T $, which are disjoint. If both occur, intersection of
the resulting members forces their single exterior points to coincide,
and each class has size at most one. If only one occurs it has size at
most two by triple capacity. Therefore $ e_T+e_{W\setminus T}\le 2$
and $|\mathcal I|\le 9+9+2=20$.
This argument subtracts the same internal occurrences before adding
external capacities. It requires a complete P9/P10 point-link and
intersection of $\mathcal I $, both already established here.
Consequently $|\F|=|N[R_i]|+|D(R_i)|\le 41+20=61$.
\end{proof}

\section{Conclusion and finite evidence boundary}
\label{sec:terminalconclusion}

\begin{proof}[Proof of the upper bound 68]
Suppose an SF4 family has at least 69 members and take a 69-member
subfamily. It remains SF4. H27 gives $\mu\ge 42$, and the universal
triple bound gives $\Delta\le 20$. Proposition~\ref{prop:f69reduction}
leaves maximum degree 16, so a degree 16 point exists.
Proposition~\ref{prop:allr16} makes its complete link R16.
Theorem~\ref{thm:r61} applies because $42\ge 41$ and $16\le 18$, and
gives $69\le 61$, a contradiction. Thus every SF4 family has size at most 68.
\end{proof}

The contradiction uses the complete local R61 theorem, including its
crowding lemma, after the nine preliminary finite blocks and the fibre
bridge. The classification and maximality of R16 supply the high-link
decompositions. The separate C5 finite lemma remains an explicit input;
the complete route and its acyclic dependency order are given in the
companion. The local R61 conclusion does not imply a global bound of $61$.

\section{The classical lower construction}
\label{sec:lower}
Partition nine points into three groups $ U_1,U_2,U_3$ of size
three, and take all quadruples formed from a pair in each of
two groups. This gives $3\cdot 3\cdot 3=27$ members, denoted
$\mathcal D_{27}$.

\begin{proposition}
$\mathcal D_{27}$ is intersecting and three-sunflower-free.
Two copies on disjoint supports give a three-sunflower-free
family of size $54$.
\end{proposition}
\begin{proof}
Two members use two of the three groups, so they share a group;
their two-element subsets in that three-point group intersect.
Suppose three distinct members formed a sunflower. A group
cannot occur in exactly two of them: the two active pairs
intersect, whereas an inactive member intersects that group
in no point. A group occurring in all three members must use
the same pair in all three. Indeed, two equal pairs and one
different pair yield unequal pairwise intersections, and the
three distinct pairs have three different singleton
intersections. Group occurrence counts are consequently
$0,1,3$. Their sum is six, since each member uses two groups.
On three groups the only possibility is $3+3+0$; the three
members then have identical pairs in the two active groups
and are equal, a contradiction. For two disjoint copies,
a triple using both copies has two intersecting members
in one copy and a third member disjoint from them. Its
pairwise intersections cannot all coincide.
\end{proof}

\section{The remaining extremal question}
\label{sec:residual}
The proof excludes $69$ members and every larger size by heredity.
It does not exclude $68$ members or establish that the classical
$54$-member construction is optimal. H27 gives only $\mu\ge41$ at
cardinality $68$, so the $\mu\ge42$ trace and fibre domains cannot be
transferred automatically. Determining $M(4,3)$ remains open.

\section{Computational evidence and reproducibility}
\label{sec:reproduction}
Each finite lemma has two obligations: a written reduction places every
candidate in the enumerated domain, and a checker verifies that entire
domain using necessary constraints. Symmetry normalization, complete
root coverage and sound monotone pruning are parts of the proof.

The companion supplies nine preliminary blocks and the final structural
checks, including C5 and both transported degree-17 cores. It contains
no optimizer dependency. Model coefficients and certificate inequalities
use integer or rational arithmetic. The full route is
\path{scripts/verify.py} with mode \texttt{complete}; mode
\texttt{quick} checks only explicit constructions and local tables.
The exact commands and finite coverage are given in \texttt{REPRODUCE.md}
and \path{proof/DEPENDENCIES.md}, relative to the repository root.
The paper and companion are available at
\url{https://github.com/aconsciousfractal/Four-Uniform-Three-Sunflower-Bound}.

The verifier requires a fresh receipt directory, checks the input manifest
and exact suite identity, and compares every scientific result with its
declared semantic outcome and coverage counters. It runs ordinary and
optimized Python and requires matching normalized results. Failure,
timeout or incomplete coverage cannot produce a complete PASS.
The verification route is exercised from an allowlisted separate copy;
a Python filesystem audit boundary restricts reads to the package,
receipts and interpreter. This is not an operating-system sandbox.

The small construction checker tests the $54$-member witness on all
$24\,804$ triples, the $45$ P8 omission pairs and $60$ P10 automorphisms.
Across five separated link models it checks $80$ near-transversal cases:
$54$ are nonempty, each with one member. A positive R16 control has seven
complete transversals, showing why the separation hypothesis is essential.

\section{Scope of verification}
H27 is imported from Proposition~6.7 of the cited preprint version;
its original exhaustive search is not re-executed here. Jin's precise
strict homomorphism theorem is stated through Nikiforov's Theorem~A.
The companion verifies the named finite computations, rather than a
proof-assistant formalization of the entire theorem. Agreement between
Python modes does not provide independence from shared mathematical
assumptions. The comparison with the published endpoint $83$ refers
specifically to the cited preprint version; no claim of priority is made.

\paragraph{Use of computational assistance.}
Generative AI tools assisted proof development, software preparation and
manuscript review. They are not authors; responsibility for the manuscript
and its claims rests with the named author.

% Inline bibliography fragment. The containing document loads hyperref.
\begin{thebibliography}{9}
\bibitem{ER1960}
P.~Erd\H{o}s and R.~Rado,
Intersection theorems for systems of sets,
\emph{Journal of the London Mathematical Society} \textbf{35} (1960), 85--90.
\href{https://doi.org/10.1112/jlms/s1-35.1.85}{doi:10.1112/jlms/s1-35.1.85}.

\bibitem{AH1974}
H.~L.~Abbott and D.~Hanson,
On finite Delta-systems,
\emph{Discrete Mathematics} \textbf{8} (1974), 1--12.
\href{https://doi.org/10.1016/0012-365X(74)90103-4}{doi:10.1016/0012-365X(74)90103-4}.

\bibitem{AHS1972}
H.~L.~Abbott, D.~Hanson and N.~Sauer,
Intersection theorems for systems of sets,
\emph{Journal of Combinatorial Theory, Series A} \textbf{12} (1972), 381--389.
\href{https://doi.org/10.1016/0097-3165(72)90103-3}{doi:10.1016/0097-3165(72)90103-3}.

\bibitem{AXB2026}
E.~Axante, C.~Budala, D.~Chitic, B.~Dumitru and M.~Nacu,
\emph{Sunflower-Free Uniform Families: Recursive Constructions and Explicit Bounds},
preprint (2026).
\href{https://arxiv.org/abs/2609.06175v1}{arXiv:2609.06175v1}.

\bibitem{JIN1993}
G.~Jin,
Triangle-Free Graphs with High Minimal Degrees,
\emph{Combinatorics, Probability and Computing} \textbf{2}(4) (1993), 479--490.
\href{https://doi.org/10.1017/S0963548300000845}{doi:10.1017/S0963548300000845}.

\bibitem{NIK2010}
V.~Nikiforov,
\emph{Chromatic number and minimum degree of \(K_r\)-free graphs},
preprint (2010).
\href{https://arxiv.org/abs/1001.2070v1}{arXiv:1001.2070v1}.

\end{thebibliography}


\end{document}
